\documentclass[a4paper,12pt]{article}
\usepackage[T2A]{fontenc}
\usepackage[utf8]{inputenc}
\usepackage[russian]{babel}
\usepackage{amsmath,amssymb,mathtools}
\usepackage{helvet}
\renewcommand{\familydefault}{\sfdefault}
\usepackage{icomma}
\usepackage[margin=21mm,headheight=14pt,headsep=8mm,footskip=12mm]{geometry}
\usepackage{microtype}
\usepackage{tikz}
\usetikzlibrary{arrows.meta,positioning,shapes.geometric,calc}
\usepackage{tabularx,booktabs,array}
\usepackage{enumitem}
\usepackage{fancyhdr}
\usepackage{setspace}
\usepackage{xcolor}
\usepackage[hidelinks]{hyperref}
\definecolor{accent}{HTML}{126E72}
\definecolor{darkaccent}{HTML}{184B50}
\definecolor{textgray}{HTML}{343B3C}
\definecolor{muted}{HTML}{657579}
\color{textgray}
\onehalfspacing
\setlength{\parindent}{0pt}
\setlength{\parskip}{5pt}
\setlist[enumerate]{leftmargin=8mm,itemsep=7pt,topsep=4pt,parsep=0pt}
\pagestyle{fancy}
\fancyhf{}
\fancyhead[L]{\footnotesize\color{muted}Финансовые задачи}
\fancyhead[R]{\footnotesize\color{muted}17.10.2026}
\fancyfoot[C]{\footnotesize\color{muted}\thepage}
\renewcommand{\headrulewidth}{0.35pt}
\renewcommand{\headrule}{\hbox to\headwidth{\color{muted}\leaders\hrule height \headrulewidth\hfill}}
\newcommand{\key}[1]{\textbf{\color{darkaccent}#1}}
\newcommand{\unit}[1]{\,\text{#1}}
\newcommand{\sechead}[1]{\vspace{7pt}{\large\bfseries\color{darkaccent}#1}\par\vspace{2pt}}
\newcolumntype{Y}{>{\raggedright\arraybackslash}X}
\renewcommand{\arraystretch}{1.35}
\begin{document}
\begin{titlepage}
\thispagestyle{empty}
\begin{center}
\vspace*{30mm}
{\large\color{muted}Математика · профильный ЕГЭ\par}
\vspace{17mm}
{\fontsize{32}{37}\selectfont\bfseries\color{darkaccent}Чек-лист\par}
\vspace{6mm}
{\fontsize{21}{28}\selectfont Финансовые задачи\par}
\vspace{4mm}
{\large Вклады и три схемы погашения кредита\par}
\vspace{20mm}
{\color{muted}17 октября 2026 года\par}
\vfill
{\normalsize\bfseries Лёвин Артём Александрович\par}
\vspace{2mm}
{\normalsize эксклюзивно для Тимофея Васильченко\par}
\vspace{20mm}
\end{center}
\begin{tikzpicture}[remember picture,overlay]
\draw[accent,line width=1.0pt]
([xshift=22mm,yshift=27mm]current page.south west)
 .. controls ([xshift=61mm,yshift=40mm]current page.south west)
 and ([xshift=-59mm,yshift=14mm]current page.south east)
 .. ([xshift=-22mm,yshift=29mm]current page.south east);
\end{tikzpicture}
\end{titlepage}

\sechead{1. Универсальный язык финансовых задач}
Обозначим через $S$ первоначальную сумму, через $r$ годовую процентную ставку, а через $p$ её десятичную запись. Удобно ввести коэффициент роста $k$:
\[
\boxed{p=\frac{r}{100},\qquad k=1+p=1+\frac{r}{100}.}
\]
Например, при $r=16\%$ имеем $p=0{,}16$, $k=1{,}16$. Если проценты начисляются раз в год, сумма $A$ после начисления становится равной $Ak$.

\key{Для кредита неизменна шапка расчётной таблицы:}
\begin{center}
\small
\begin{tabularx}{\linewidth}{@{}cYYYYY@{}}
\toprule
Год & Начальная сумма & Начисленный процент & Общая задолженность & Платёж & Остаток\\
\midrule
$i$ & $B_i$ & $B_ip$ & $B_ik$ & $X_i$ & $B_ik-X_i$\\
\bottomrule
\end{tabularx}
\end{center}
\key{Связь строк:} остаток после платежа в предыдущем году становится начальной суммой следующего. В конце срока кредита остаток равен нулю.

\sechead{2. Вклад с дополнительным взносом}
Если в начале второго года внесены дополнительные $D$ рублей, а проценты начисляются в конце каждого года, то:
\[
\underbrace{S\vphantom{(}}_{\text{вклад}}
\;\xrightarrow{\times k}\;Sk
\;\xrightarrow{+D}\;Sk+D
\;\xrightarrow{\times k}\;\boxed{(Sk+D)k}.
\]
\key{Пример занятия.} В январе 2031 года внесли $850$ тыс. руб., в январе 2032 года добавили $350$ тыс. руб. К январю 2033 года на счёте $1549{,}76$ тыс. руб. Найдём $r$.
\[
(850k+350)k=1549{,}76,
\quad 850k^2+350k-1549{,}76=0.
\]
Положительный корень $k=1{,}16$; второй корень отрицателен и не подходит, поскольку $k>0$. Следовательно, $r=(k-1)\cdot100=\boxed{16\%}$.

\key{Важно.} Дополнительный взнос включают в сумму \emph{до} начисления процентов за следующий год. Процентную ставку и коэффициент $k$ не следует путать.

\clearpage
\sechead{3. Как узнать схему кредитования по условию}
\begin{tabularx}{\linewidth}{@{}p{45mm}Y Y@{}}
\toprule
Схема & Ключевые слова в условии & Что остаётся одинаковым\\
\midrule
\textbf{Аннуитетная} & «равными платежами» & $X_1=X_2=\cdots=X_n=x$\\
\textbf{Дифференцированная} & «долг уменьшается на одну и ту же сумму» & Доля погашения \emph{основного долга} $S/n$\\
\textbf{Произвольная} & «долг становится на $a\%$ меньше долга предыдущего июля» & Один и тот же \emph{процент уменьшения остатка}\\
\bottomrule
\end{tabularx}

\key{Ловушка.} Слова «на $20\%$ меньше предыдущего долга» не означают равные платежи; они описывают уменьшение остатка в определённой пропорции.

\sechead{4. Произвольная схема: задан процент уменьшения долга}
\key{Пример занятия.} Кредит $S=1{,}5$ млн руб. выдан на четыре года под $10\%$ годовых. В июле 2031, 2032 и 2033 годов долг должен становиться на $40\%$ меньше июльского долга предыдущего года. К июлю 2034 года кредит погашен.

Положим $q=1-0{,}40=0{,}60$, $p=0{,}10$. Остатки в конце первых трёх периодов равны $qS$, $q^2S$, $q^3S$. Каждый платёж найдём как \emph{общую задолженность минус остаток}.
\begin{center}
\small
\begin{tabularx}{\linewidth}{@{}cYYYYY@{}}
\toprule
Год & Начальная сумма & Процент & Общий долг & Платёж & Остаток\\
\midrule
1 & $S$ & $Sp$ & $S(1+p)$ & $S(1+p)-0{,}6S$ & $0{,}6S$\\
2 & $0{,}6S$ & $0{,}6Sp$ & $0{,}6S(1+p)$ & $0{,}6S(1+p)-0{,}36S$ & $0{,}36S$\\
3 & $0{,}36S$ & $0{,}36Sp$ & $0{,}36S(1+p)$ & $0{,}36S(1+p)-0{,}216S$ & $0{,}216S$\\
4 & $0{,}216S$ & $0{,}216Sp$ & $0{,}216S(1+p)$ & $0{,}216S(1+p)$ & $0$\\
\bottomrule
\end{tabularx}
\end{center}
Искомый платёж четвёртого года:
\[
X_4=0{,}216\cdot1{,}5\cdot1{,}10
=0{,}3564\unit{млн руб.}=\boxed{356400\unit{руб.}}
\]
\key{Обобщение:} если в первые $n-1$ лет долг каждый раз уменьшается до доли $q$ от прежнего, то окончательный платёж
\[
\boxed{X_n=Sq^{n-1}(1+p)}.
\]

\clearpage
\sechead{5. Аннуитет: равные платежи}
Пусть $x$ --- одинаковый платёж в конце каждого из $n$ расчётных периодов. Для двух лет:
\begin{center}
\small
\begin{tabularx}{\linewidth}{@{}cYYYYY@{}}
\toprule
Год & Начальная сумма & Процент & Общий долг & Платёж & Остаток\\
\midrule
1 & $S$ & $Sp$ & $Sk$ & $x$ & $Sk-x$\\
2 & $Sk-x$ & $(Sk-x)p$ & $(Sk-x)k$ & $x$ & $(Sk-x)k-x=0$\\
\bottomrule
\end{tabularx}
\end{center}
Отсюда $Sk^2-x(k+1)=0$, поэтому
\[
\boxed{x=\frac{Sk^2}{k+1}},\qquad \text{общая сумма платежей}=2x.
\]
\key{Пример занятия.} $S=810$ тыс. руб., $r=16\%$, $k=1{,}16$.
\[
x=\frac{810\cdot1{,}16^2}{2{,}16}=504{,}6\unit{тыс. руб.},
\quad 2x=\boxed{1009200\unit{руб.}}
\]
Для трёх лет аналогично: $Sk^3-x(k^2+k+1)=0$, откуда
\[
\boxed{x=\frac{Sk^3}{k^2+k+1}},\qquad \text{общая сумма}=3x.
\]

\sechead{6. Дифференцированная схема}
Если основной долг уменьшается на одну и ту же сумму в каждом из $n$ периодов, после $i$-го платежа остаётся $S(1-i/n)$. Поэтому
\[
\boxed{X_i=\frac Sn+pS\left(1-\frac{i-1}{n}\right)},
\qquad i=1,2,\ldots,n.
\]
Здесь $S/n$ --- постоянная доля основного долга, а процентная часть платежа уменьшается по мере погашения кредита. \key{Платежи при этой схеме, как правило, различаются.}

\clearpage
\thispagestyle{fancy}
\begin{center}
{\Large\bfseries\color{darkaccent}Алгоритм: как выбрать способ решения}\par
\end{center}
\vspace{13mm}
\begin{center}
\begin{tikzpicture}[
start/.style={ellipse,draw=accent,line width=.9pt,minimum height=12mm,text width=39mm,align=center,font=\small},
question/.style={diamond,draw=darkaccent,line width=.85pt,aspect=2.4,text width=44mm,inner sep=1pt,align=center,font=\small},
result/.style={rectangle,rounded corners=2pt,draw=accent,line width=.9pt,text width=52mm,minimum height=17mm,align=center,font=\small},
arrow/.style={-Latex,thick,draw=darkaccent},
branch/.style={font=\footnotesize,fill=white,inner sep=2pt}
]
\node[start] (st) at (-3.3,1.1) {Прочитайте условие};
\node[question] (q1) at (-3.3,-1.6) {Это вклад с начислением процентов?};
\node[result] (a1) at (3.3,-1.6) {\textbf{Вклад}\vspace{1mm}\\Отмечайте взносы перед соответствующим начислением процентов};
\node[question] (q2) at (-3.3,-5.2) {Кредит погашают равными платежами?};
\node[result] (a2) at (3.3,-5.2) {\textbf{Аннуитетная схема}\vspace{1mm}\\Каждый платёж $x$; последний остаток равен нулю};
\node[question] (q3) at (-3.3,-8.8) {Основной долг уменьшается на одну и ту же сумму?};
\node[result] (a3) at (3.3,-8.8) {\textbf{Дифференцированная схема}\vspace{1mm}\\Основной долг уменьшается на $S/n$ за период};
\node[result] (a4) at (-3.3,-12.2) {\textbf{Произвольная схема}\vspace{1mm}\\Запишите требуемые остатки, затем найдите платежи по таблице};
\draw[arrow] (st) -- (q1);
\draw[arrow] (q1.east) -- node[branch,above]{да} (a1.west);
\draw[arrow] (q1.south) -- node[branch,right]{нет} (q2.north);
\draw[arrow] (q2.east) -- node[branch,above]{да} (a2.west);
\draw[arrow] (q2.south) -- node[branch,right]{нет} (q3.north);
\draw[arrow] (q3.east) -- node[branch,above]{да} (a3.west);
\draw[arrow] (q3.south) -- node[branch,right]{нет} (a4.north);
\end{tikzpicture}
\end{center}
\clearpage

\sechead{7. Тренировка --- Путь А}
\textit{Во всех задачах проценты начисляются один раз в год до очередного платежа или в конце очередного года вклада. Составляйте таблицу там, где это помогает обосновать решение. Вычисления выполняйте без промежуточного округления.}
\begin{enumerate}[label=\textbf{\arabic*.}]
\item На вклад в размере $400$ тыс. руб. за год начислили $12\%$. Найдите коэффициент роста $k$ и сумму на счёте в конце года.
\item В начале первого года открыли вклад на $500$ тыс. руб. под $10\%$ годовых. В начале второго года дополнительно внесли $200$ тыс. руб. Какова сумма на счёте после начисления процентов за второй год?
\item В начале первого года на вклад внесли $600$ тыс. руб.; в начале второго года добавили $250$ тыс. руб. Через два года на счёте оказалось $1164$ тыс. руб. Найдите годовую ставку $r>0$.
\item Кредит $800$ тыс. руб. взяли на два года под $10\%$ годовых. В конце первого года остаток долга должен стать на $25\%$ меньше первоначальной суммы, а в конце второго года кредит будет полностью погашен. Найдите второй платёж.
\item Кредит $1$ млн руб. взяли на три года под $12\%$ годовых. В конце первого и второго годов остаток долга каждый раз становится на $20\%$ меньше остатка предыдущего года; в конце третьего года долг равен нулю. Найдите третий платёж.
\item Кредит $1{,}25$ млн руб. взяли на четыре года под $8\%$ годовых. В конце каждого из первых трёх годов остаток долга уменьшается на $40\%$ по сравнению с предыдущим; в четвёртом году долг погашается полностью. Найдите четвёртый платёж.
\item Кредит $900$ тыс. руб. взяли на три года под $10\%$ годовых. После каждого ежегодного платежа основной долг уменьшается на одинаковую сумму. Найдите второй платёж.
\item Кредит $1{,}2$ млн руб. взяли на четыре года под $15\%$ годовых. После каждого ежегодного платежа основной долг уменьшается на одинаковую сумму. Найдите общую сумму всех платежей.
\item Кредит $210$ тыс. руб. взяли на два года под $10\%$ годовых и погасили двумя равными ежегодными платежами. Найдите общую сумму платежей.
\item Кредит $331$ тыс. руб. взяли на три года под $10\%$ годовых и погасили тремя равными ежегодными платежами. Найдите величину каждого платежа и их общую сумму.
\end{enumerate}

\clearpage
\sechead{8. Ответы для самопроверки}
\begin{center}
\renewcommand{\arraystretch}{1.55}
\begin{tabularx}{\linewidth}{@{}cY@{}}
\toprule
\textbf{№} & \textbf{Ответ}\\
\midrule
1 & $k=1{,}12$; $448\unit{тыс. руб.}$\\
2 & $825\unit{тыс. руб.}$\\
3 & $20\%$\\
4 & $660000\unit{руб.}$\\
5 & $716800\unit{руб.}$\\
6 & $291600\unit{руб.}$\\
7 & $360000\unit{руб.}$\\
8 & $1650000\unit{руб.}$\\
9 & $242000\unit{руб.}$\\
10 & Каждый платёж $133100\unit{руб.}$; общая сумма $399300\unit{руб.}$\\
\bottomrule
\end{tabularx}
\end{center}
\vspace{10mm}
\key{Контроль решения.} Проверьте, что проценты начислялись на фактический остаток долга, число строк соответствует сроку кредита, а в последней строке остаток равен нулю. Если вопрос об \emph{общей сумме платежей}, сложите все платежи.
\end{document}
